\documentclass[11pt,a4paper]{article}
\usepackage{fontspec}
\usepackage{kotex}
\usepackage{geometry}
\usepackage{graphicx}
\usepackage{xcolor}
\usepackage{tcolorbox}
\usepackage{booktabs}
\usepackage{array}
\usepackage{longtable}
\usepackage{colortbl}
\usepackage{fancyhdr}
\usepackage{titlesec}
\usepackage{hyperref}
\usepackage{emoji}
\usepackage{amsmath}
\usepackage{amssymb}
\usepackage{enumitem}

\geometry{left=25mm, right=25mm, top=30mm, bottom=30mm}

\setmainfont{Noto Sans CJK KR}
\setsansfont{Noto Sans CJK KR}
\setmonofont{Noto Sans CJK KR}
\setemojifont{Noto Color Emoji}

\definecolor{primary}{HTML}{1976D2}
\definecolor{darkbrown}{HTML}{3E2723}
\definecolor{mediumbrown}{HTML}{5D4037}
\definecolor{lightbrown}{HTML}{8D6E63}
\definecolor{cream}{HTML}{FFF3E0}
\definecolor{lightgray}{HTML}{F5F5F5}
\definecolor{tableheader}{HTML}{E8E8E8}

\titleformat{\section}
  {\Large\bfseries\color{primary}}
  {\thesection.}{0.5em}{}
\titleformat{\subsection}
  {\large\bfseries\color{darkbrown}}
  {\thesubsection}{0.5em}{}
\titleformat{\subsubsection}
  {\normalsize\bfseries\color{mediumbrown}}
  {\thesubsubsection}{0.5em}{}

\renewcommand{\contentsname}{목차}
\renewcommand{\figurename}{그림}
\renewcommand{\tablename}{표}

\tcbset{
  tipbox/.style={colback=cream,colframe=primary,boxrule=1.5pt,arc=3mm,
    left=5mm,right=5mm,top=3mm,bottom=3mm,fonttitle=\bfseries\color{primary}},
  formulabox/.style={colback=lightgray,colframe=lightgray,boxrule=0pt,arc=2mm,
    left=5mm,right=5mm,top=3mm,bottom=3mm},
  summarybox/.style={colback=white,colframe=primary,boxrule=2pt,arc=4mm,
    left=5mm,right=5mm,top=5mm,bottom=5mm}
}

\hypersetup{colorlinks=true,linkcolor=primary,urlcolor=primary,bookmarks=true,unicode=true}


\pagestyle{fancy}
\fancyhf{}
\fancyhead[R]{\small\color{lightbrown}마이크로비트 Grove I2C 센서 시뮬레이터 퀵가이드 v1.0.0.0}
\fancyfoot[C]{\thepage}
\renewcommand{\headrulewidth}{0.4pt}

\begin{document}

\begin{titlepage}
\centering
\vspace*{3cm}
{\Huge\bfseries\color{primary}\emoji{electric-plug} 마이크로비트 Grove I2C 센서 시뮬레이터}\\[0.5cm]
{\Huge\bfseries\color{primary} 퀵가이드}
\\[1.5cm]
{\Large 시뮬레이터 버전: v9.2.0.0}
\\[2cm]
\begin{tcolorbox}[summarybox]
\centering
{\large 이 퀵가이드는 마이크로비트 Grove I2C 센서 시뮬레이터와 함께 보는 내비게이션입니다.}
\end{tcolorbox}
\vfill
{\large 사물인터넷과 빅데이터}\\[0.3cm]
{\large 청주교육대학교 컴퓨터교육과}
\vspace{1cm}
\end{titlepage}

\section*{\emoji{open-book} 이 문서 사용법}

\begin{tcolorbox}[tipbox]
\textbf{이 퀵가이드는 센서 시뮬레이터와 함께 보는 내비게이션입니다.}

\end{tcolorbox}

\begin{center}
\renewcommand{\arraystretch}{1.4}
\begin{tabular}{|p{8.2cm}|p{6.8cm}|}
\hline
\rowcolor{tableheader}
\textbf{상황} & \textbf{참조 섹션} \\
\hline
처음 시작 & \emoji{high-voltage} 초속성 실습 \\
\hline
어떤 센서가 있는지 모름 & \emoji{magnifying-glass-tilted-left} 센서 한눈에 보기 \\
\hline
버튼·영역 위치 모름 & \emoji{desktop-computer} 화면 구성 \\
\hline
문제 발생 & \emoji{sos} 트러블슈팅 \\
\hline
\end{tabular}
\end{center}


\section*{\emoji{high-voltage} 1분 초속성 실습}

\begin{center}
\renewcommand{\arraystretch}{1.4}
\begin{tabular}{|p{1.5cm}|p{8.9cm}|p{4.6cm}|}
\hline
\rowcolor{tableheader}
\textbf{단계} & \textbf{행동} & \textbf{확인} \\
\hline
1 & 오른쪽 센서 목록에서 \textbf{SHT40 온습도} 카드를 찾으세요 & \emoji{thermometer} 아이콘 보이나요? \\
\hline
2 & 카드를 \textbf{드래그}하여 중앙 포트 \texttt{[1]}에 \textbf{드롭}하세요 & 포트가 파란색으로 바뀌었나요? \\
\hline
3 & 하단 \textbf{\emoji{bar-chart} 측정 데이터}에서 "정밀 온도"와 "습도"를 확인하세요 & 데이터 항목이 추가되었으면 \emoji{party-popper} \\
\hline
\end{tabular}
\end{center}


\section*{\emoji{rocket} 빠른 시작 (3단계)}

\subsection*{\emoji{keycap-1} 센서 탐색하기}

오른쪽 \textbf{\emoji{electric-plug} 외장 Grove 센서} 영역에 10개의 센서 카드가 있습니다. 카드를 \textbf{마우스 오버}하면 센서 상세 설명이 툴팁으로 나타납니다.


무엇을 측정할지 모르겠다면? \textbf{\emoji{clipboard} 추천 키워드} 버튼을 클릭하세요. "온도 습도", "밝기 조도", "거리 장애물" 등 13가지 카테고리에서 원하는 키워드를 선택하면 관련 센서가 \textbf{하이라이트}됩니다.


\subsection*{\emoji{keycap-2} 센서 연결하기}

하이라이트된 센서 카드를 \textbf{드래그}하여 중앙 Grove Shield의 빈 포트(1\textasciitilde{}6)에 \textbf{드롭}합니다. 포트에 센서가 연결되면 포트 색상이 변하고, 센서 이름이 포트에 표시됩니다.


최대 \textbf{6개 센서}를 동시에 연결할 수 있습니다. 연결된 센서를 다른 포트로 옮기려면 포트의 센서를 드래그하세요.


\subsection*{\emoji{keycap-3} 측정 데이터 확인하기}

화면 하단 \textbf{\emoji{bar-chart} 측정 데이터} 패널에서 현재 연결된 센서들의 데이터 항목을 확인합니다. 내장 센서 7항목은 항상 표시되고, 외장 센서 데이터는 연결할 때마다 추가됩니다.



\section*{\emoji{magnifying-glass-tilted-left} 센서 한눈에 보기}

\subsection*{내장 센서 (마이크로비트 V2 — 항상 사용 가능)}

\begin{center}
\renewcommand{\arraystretch}{1.4}
\begin{tabular}{|p{3.6cm}|p{4.5cm}|p{6.9cm}|}
\hline
\rowcolor{tableheader}
\textbf{센서} & \textbf{측정} & \textbf{프로젝트 예시} \\
\hline
\emoji{person-running} 가속도계 & X/Y/Z축 가속도 & 걸음수, 기울기, 진동 \\
\hline
\emoji{compass} 나침반 & 방위각 (0\textasciitilde{}360°) & 방향 탐지 \\
\hline
\emoji{thermometer} 온도 & 대략적 온도 & 교실 환경 (정밀도 낮음) \\
\hline
\emoji{light-bulb} 빛 & 밝기 수준 & 조도 변화, 그림자 \\
\hline
\emoji{microphone} 마이크 & 소리 크기 & 소음 측정, 박수 감지 \\
\hline
\emoji{backhand-index-pointing-up} 터치 로고 & 터치 감지 & 버튼 입력 \\
\hline
\end{tabular}
\end{center}

\subsection*{외장 센서 (Grove I2C — 포트에 연결해서 사용)}

\begin{center}
\renewcommand{\arraystretch}{1.4}
\begin{tabular}{|p{5.3cm}|p{4.4cm}|p{5.3cm}|}
\hline
\rowcolor{tableheader}
\textbf{센서} & \textbf{측정} & \textbf{프로젝트 예시} \\
\hline
\emoji{thermometer} SHT40 온습도 & 정밀 온도·습도 & 교실 쾌적도 (정밀) \\
\hline
\emoji{sun-behind-small-cloud} BMP280 기압 & 기압·고도 & 날씨 예보 \\
\hline
\emoji{bright-button} TSL2561 조도 & 정밀 조도 (lux) & 조명 환경 \\
\hline
\emoji{artist-palette} TCS34725 컬러 & RGB 색상값 & 색상 분류 \\
\hline
\emoji{wind-face} SGP30 공기질 & CO\textsubscript{2}·TVOC & 환기 시점 \\
\hline
\emoji{triangular-ruler} LIS3DHTR 가속도 & 정밀 3축 가속도 & 진동, 충격 \\
\hline
\emoji{straight-ruler} VL53L0X 거리 & ToF 거리 (mm) & 장애물 감지 \\
\hline
\emoji{musical-keyboard} MPR121 터치 & 12채널 터치 & 터치 피아노 \\
\hline
\emoji{waving-hand} PAJ7620 제스처 & 9방향 제스처 & 손동작 인식 \\
\hline
\emoji{television} OLED 디스플레이 & 화면 출력 & 측정값 표시 \\
\hline
\end{tabular}
\end{center}

\begin{tcolorbox}[summarybox]
\emoji{light-bulb} \textbf{팁}: 내장 온도보다 SHT40이 훨씬 정밀합니다. 온도 관련 프로젝트라면 SHT40을 추천합니다!

\end{tcolorbox}


\section*{\emoji{desktop-computer} 화면 구성}

\begin{center}
\includegraphics[width=0.95\textwidth]{images/S4_GroveSensor_screen.png}
\end{center}

\begin{center}
\renewcommand{\arraystretch}{1.4}
\begin{tabular}{|p{3.5cm}|p{2.6cm}|p{8.9cm}|}
\hline
\rowcolor{tableheader}
\textbf{영역} & \textbf{위치} & \textbf{역할} \\
\hline
내장 센서 & 좌측 상단 & 마이크로비트 기본 탑재 센서 6종 \\
\hline
마이크로비트 & 좌측 중앙 & SVG로 그린 마이크로비트 보드 \\
\hline
데이터로거 & 좌측 하단 & RTC+SD카드 (항상 장착) \\
\hline
Grove Shield & 중앙 & I2C 포트 6개 — 센서 드롭 영역 \\
\hline
외장 센서 & 우측 & Grove I2C 센서 10종 카드 \\
\hline
검색 & 우측 하단 & 키워드 검색 + 추천 키워드 \\
\hline
측정 데이터 & 하단 전체 & 내장 + 외장 데이터 항목 목록 \\
\hline
\end{tabular}
\end{center}


\section*{\emoji{bullseye} 관찰 포인트}

\begin{center}
\renewcommand{\arraystretch}{1.4}
\begin{tabular}{|p{3.3cm}|p{4.5cm}|p{7.2cm}|}
\hline
\rowcolor{tableheader}
\textbf{관찰 대상} & \textbf{찾아볼 장면} & \textbf{의미} \\
\hline
포트 색상 변화 & 센서 드롭 시 포트가 파란색으로 & 연결 성공 \\
\hline
데이터 패널 갱신 & 센서 연결 시 하단에 항목 추가 & "이 센서로 이런 데이터를 측정할 수 있구나" \\
\hline
센서 카드 하이라이트 & 검색 시 관련 카드 테두리 강조 & 내 프로젝트에 맞는 센서 발견 \\
\hline
내장 vs 외장 배지 색상 & \emoji{green-circle}초록 = 내장, \emoji{blue-circle}파랑 = 외장 & 추가 장비 필요 여부 판단 \\
\hline
데이터 항목 수 증가 & 센서 추가할수록 데이터 행 증가 & "센서가 많으면 더 다양한 데이터를 모을 수 있다" \\
\hline
\end{tabular}
\end{center}


\section*{\emoji{sos} 도와주세요!}

\begin{center}
\renewcommand{\arraystretch}{1.4}
\begin{tabular}{|p{3.8cm}|p{5.2cm}|p{6.0cm}|}
\hline
\rowcolor{tableheader}
\textbf{증상} & \textbf{원인} & \textbf{해결책} \\
\hline
센서를 드래그할 수 없어요 & 외장 센서만 드래그 가능 & 오른쪽 \emoji{blue-circle}파랑 테두리 카드를 드래그하세요 \\
\hline
포트에 드롭이 안 돼요 & 이미 같은 센서가 다른 포트에 연결됨 & 기존 연결을 해제 후 다시 시도 \\
\hline
검색 결과가 없어요 & 키워드가 너무 구체적 & \textbf{\emoji{clipboard} 추천 키워드}에서 선택하세요 \\
\hline
데이터 패널이 비어있어요 & 아직 외장 센서 미연결 & 센서 1개 이상 연결하면 외장 데이터 표시 \\
\hline
화면이 깨져요 & 브라우저 창이 너무 작음 & 창을 넓히거나, 확대/축소(Ctrl+0) 초기화 \\
\hline
\end{tabular}
\end{center}


\section*{\emoji{key} 핵심 용어}

\begin{center}
\renewcommand{\arraystretch}{1.4}
\begin{tabular}{|p{1.5cm}|p{3.2cm}|p{4.4cm}|p{5.9cm}|}
\hline
\rowcolor{tableheader}
\textbf{용어} & \textbf{학술 용어} & \textbf{의미} & \textbf{쉬운 설명} \\
\hline
I2C & Inter-Integrated Circuit & 2선식 직렬 통신 규격 & 센서를 연결하는 전선 2개짜리 규격 \\
\hline
Grove & Grove System & Seeed Studio의 커넥터 규격 & 레고처럼 꽂기만 하면 되는 센서 연결 방식 \\
\hline
I2C 주소 & I2C Address & 각 센서의 고유 번호 & 버스에서 센서를 구별하는 이름표 (예: 0x44) \\
\hline
내장 센서 & Built-in Sensor & 마이크로비트에 기본 탑재된 센서 & 추가 장비 없이 바로 쓸 수 있는 센서 \\
\hline
외장 센서 & External Sensor & Grove Shield로 추가 연결하는 센서 & 더 정밀하거나 특별한 측정을 위해 추가하는 센서 \\
\hline
데이터로거 & Data Logger & 데이터 기록 장치 & 센서 값을 시간표와 함께 파일로 저장하는 장치 \\
\hline
\end{tabular}
\end{center}


\section*{\emoji{thought-balloon} 생각해보기}

\begin{tcolorbox}[summarybox]
\textbf{내 프로젝트에 센서가 몇 개 필요한가?}

\vspace{0.3em}
내장 센서만으로 충분한가요, 외장 센서가 꼭 필요한가요?

"측정할 수 있는 것"과 "측정하고 싶은 것"의 차이를 생각해보세요.

센서가 많다고 좋은 프로젝트가 되는 건 아닙니다!

\end{tcolorbox}


\section*{참고자료}

\begin{center}
\renewcommand{\arraystretch}{1.4}
\begin{tabular}{|p{1.5cm}|p{8.3cm}|p{2.4cm}|p{2.8cm}|}
\hline
\rowcolor{tableheader}
\textbf{\#} & \textbf{자료명} & \textbf{유형} & \textbf{비고} \\
\hline
1 & 마이크로비트Grove센서시뮬레이터\_v9.2.0.0 & HTML 파일 & 시뮬레이터 \\
\hline
2 & 마이크로비트Grove센서시뮬레이터\_개발참조\_v1.0.0.0 & 문서 & 기술 상세 \\
\hline
3 & 코딩관리지침\_v4.0.0.0 \S{}17 & 프로젝트 문서 & 퀵가이드 작성법 \\
\hline
\end{tabular}
\end{center}


\vfill
\begin{center}
\small\color{lightbrown}
\textit{마이크로비트 Grove I2C 센서 시뮬레이터 퀵가이드}
\end{center}
\end{document}